\documentclass[11pt,a4paper]{article}
\input{../template/utm-lab}

\newcommand{\labnumber}{12}
\newcommand{\labtitle}{Producer, Consumer and the Other Tools}
\newcommand{\studentname}{Daniel Nicorici}
\newcommand{\studentgroup}{FAF-243}

\begin{document}

\utmtitlepage
\tableofcontents
\newpage

\plainsection{Purpose of the work}

In Lab 11 one semaphore was used at a time: as a lock, as a signal, as a counter. The purpose of
this lab is to make the tools work together: to solve the two classic problems, producer and
consumer, readers and writers, to meet the condition variable, and to build a deadlock and then
remove it.

The work was done on my Arch Linux laptop with gcc 16.2.1, in the folder
\code{\textasciitilde/os-lab12}. All programs were compiled with \code{-pthread}.

\section{Part 1. Producer and consumer}

A producer puts items into a buffer of 4 slots and a slow consumer takes them out. Three tools
work together: \code{empty} counts the free slots and starts at 4, \code{full} counts the filled
slots and starts at 0, and a mutex protects the buffer itself.

\begin{lstlisting}[style=ccode,caption={prodcons.c}]
#include <stdio.h>
#include <unistd.h>
#include <pthread.h>
#include <semaphore.h>

#define SIZE 4
#define ITEMS 8

int buf[SIZE];
int in = 0, out = 0, count = 0;
sem_t empty, full;
pthread_mutex_t lock = PTHREAD_MUTEX_INITIALIZER;

void *producer(void *arg) {
    for (int item = 1; item <= ITEMS; item++) {
        sem_wait(&empty);
        pthread_mutex_lock(&lock);
        buf[in] = item;
        in = (in + 1) % SIZE;
        count++;
        printf("put %d    (in buffer: %d)\n", item, count);
        pthread_mutex_unlock(&lock);
        sem_post(&full);
    }
    return NULL;
}

void *consumer(void *arg) {
    for (int k = 0; k < ITEMS; k++) {
        sem_wait(&full);
        pthread_mutex_lock(&lock);
        int item = buf[out];
        out = (out + 1) % SIZE;
        count--;
        printf("            take %d   (in buffer: %d)\n", item, count);
        pthread_mutex_unlock(&lock);
        sem_post(&empty);
        usleep(300000);
    }
    return NULL;
}

int main(void) {
    sem_init(&empty, 0, SIZE);
    sem_init(&full, 0, 0);
    pthread_t p, c;
    pthread_create(&c, NULL, consumer, NULL);
    pthread_create(&p, NULL, producer, NULL);
    pthread_join(p, NULL);
    pthread_join(c, NULL);
    sem_destroy(&empty);
    sem_destroy(&full);
    return 0;
}
\end{lstlisting}

\begin{lstlisting}[caption={prodcons}]
$ gcc -Wall -pthread -o prodcons prodcons.c
$ ./prodcons
put 1    (in buffer: 1)
put 2    (in buffer: 2)
            take 1   (in buffer: 1)
put 3    (in buffer: 2)
put 4    (in buffer: 3)
put 5    (in buffer: 4)
            take 2   (in buffer: 3)
put 6    (in buffer: 4)
            take 3   (in buffer: 3)
put 7    (in buffer: 4)
            take 4   (in buffer: 3)
put 8    (in buffer: 4)
            take 5   (in buffer: 3)
            take 6   (in buffer: 2)
            take 7   (in buffer: 1)
            take 8   (in buffer: 0)
\end{lstlisting}

Then, in \code{producer()}, the mutex was locked before \code{sem\_wait(\&empty)}:

\begin{lstlisting}[style=ccode,caption={The two lines of producer(), swapped}]
        pthread_mutex_lock(&lock);
        sem_wait(&empty);
\end{lstlisting}

\begin{lstlisting}[caption={prodcons with the lines swapped}]
$ gcc -Wall -pthread -o prodcons prodcons.c && timeout 5 ./prodcons; echo "exit code: $?"
put 1    (in buffer: 1)
put 2    (in buffer: 2)
put 3    (in buffer: 3)
put 4    (in buffer: 4)
exit code: 124
\end{lstlisting}

\textbf{Observations:}
\begin{enumerate}
  \item \textbf{What stops the fast producer.} \code{sem\_wait(\&empty)}. \code{empty} starts
  at 4 and loses one unit for every item put in; at 0 the producer blocks until the consumer
  takes an item and calls \code{sem\_post(\&empty)}. That is why ``in buffer'' never goes above
  4.
  \item \textbf{Who waits at the start.} The consumer, at \code{sem\_wait(\&full)}: \code{full}
  starts at 0 because the buffer is empty, so there is nothing to take until the producer posts.
  \item \textbf{The swapped lines.} The program puts four items and then hangs; \code{timeout}
  had to kill it (exit code 124). When the buffer is full, the producer blocks in
  \code{sem\_wait(\&empty)} while it still holds the mutex. The consumer is the only one that
  can free a slot, but it needs the same mutex to take an item. Each waits for the other
  forever: a deadlock. In the correct order a thread never sleeps while holding the mutex. After
  the test the lines were put back.
\end{enumerate}

\section{Part 2. The condition variable}

The problem of \code{order.c} from Lab 11, solved another way. The state is a normal variable,
\code{ready}, protected by a mutex. The condition variable is only the place where B sleeps
until A says that something changed. \code{pthread\_cond\_wait} releases the mutex, sleeps, and
takes the mutex again before it returns.

\begin{lstlisting}[style=ccode,caption={condvar.c}]
#include <stdio.h>
#include <string.h>
#include <unistd.h>
#include <pthread.h>

char message[32];
int ready = 0;
pthread_mutex_t lock = PTHREAD_MUTEX_INITIALIZER;
pthread_cond_t changed = PTHREAD_COND_INITIALIZER;

void *writer(void *arg) {
    sleep(1);
    pthread_mutex_lock(&lock);
    strcpy(message, "hello");
    ready = 1;
    printf("A: message written\n");
    pthread_cond_signal(&changed);
    pthread_mutex_unlock(&lock);
    return NULL;
}

void *reader(void *arg) {
    pthread_mutex_lock(&lock);
    while (!ready) {
        printf("B: not ready, I sleep\n");
        pthread_cond_wait(&changed, &lock);
    }
    printf("B: I read \"%s\"\n", message);
    pthread_mutex_unlock(&lock);
    return NULL;
}

int main(void) {
    pthread_t a, b;
    pthread_create(&b, NULL, reader, NULL);
    pthread_create(&a, NULL, writer, NULL);
    pthread_join(a, NULL);
    pthread_join(b, NULL);
    return 0;
}
\end{lstlisting}

\begin{lstlisting}[caption={condvar, with while and with if}]
$ gcc -Wall -pthread -o condvar condvar.c
$ ./condvar
B: not ready, I sleep
A: message written
B: I read "hello"
$ gcc -Wall -pthread -o condvar condvar.c && ./condvar
B: not ready, I sleep
A: message written
B: I read "hello"
\end{lstlisting}

The second run is the version with \code{if (!ready)} instead of \code{while (!ready)}.

\textbf{Observations:}
\begin{enumerate}
  \item \textbf{Which variable remembers.} The ordinary variable \code{ready}. A condition
  variable has no memory: a signal sent when nobody is waiting is lost, so the fact itself must
  be stored in a variable protected by the mutex.
  \item \textbf{With \code{if}.} It still works here and prints the same three lines, because
  there is one reader, one writer and one signal.
  \item \textbf{Why \code{while}.} After waking, the thread must check the condition again. It
  can wake and still find \code{ready == 0} because of a spurious wakeup, which
  \code{pthread\_cond\_wait} is allowed to produce, or because another waiting thread woke
  first and already used up what was ready.
\end{enumerate}

\section{Part 3. Readers and writers}

Many readers may read together; a writer must be alone. Two semaphores are used: \code{room},
where 1 means free for a writer, and \code{rlock}, a lock around the counter \code{readers}.
The first reader in locks the room, the last one out unlocks it.

\begin{lstlisting}[style=ccode,caption={readers.c}]
#include <stdio.h>
#include <unistd.h>
#include <pthread.h>
#include <semaphore.h>

#define READERS 3
#define WRITERS 2

int value = 0;
int readers = 0;
sem_t rlock;
sem_t room;

void *reader(void *arg) {
    long id = (long)arg;
    sem_wait(&rlock);
    if (++readers == 1)
        sem_wait(&room);
    sem_post(&rlock);
    printf("reader %ld reads %d     (readers inside: %d)\n", id, value, readers);
    usleep(200000);
    sem_wait(&rlock);
    if (--readers == 0)
        sem_post(&room);
    sem_post(&rlock);
    return NULL;
}

void *writer(void *arg) {
    long id = (long)arg;
    sem_wait(&room);
    value += 10;
    printf("writer %ld writes %d (alone)\n", id, value);
    usleep(200000);
    sem_post(&room);
    return NULL;
}

int main(void) {
    sem_init(&rlock, 0, 1);
    sem_init(&room, 0, 1);
    pthread_t r[READERS], w[WRITERS];
    pthread_create(&w[0], NULL, writer, (void *)1L);
    for (long i = 0; i < READERS; i++)
        pthread_create(&r[i], NULL, reader, (void *)(i + 1));
    pthread_create(&w[1], NULL, writer, (void *)2L);
    for (int i = 0; i < READERS; i++) pthread_join(r[i], NULL);
    for (int i = 0; i < WRITERS; i++) pthread_join(w[i], NULL);
    return 0;
}
\end{lstlisting}

\begin{lstlisting}[caption={readers}]
$ gcc -Wall -pthread -o readers readers.c
$ ./readers
writer 1 writes 10 (alone)
reader 1 reads 10     (readers inside: 1)
reader 2 reads 10     (readers inside: 2)
reader 3 reads 10     (readers inside: 3)
writer 2 writes 20 (alone)
\end{lstlisting}

\textbf{Observations:}
\begin{enumerate}
  \item \textbf{The two lines.} \code{if (++readers == 1) sem\_wait(\&room);} and
  \code{if (--readers == 0) sem\_post(\&room);}. Only the first reader takes the room and only
  the last one gives it back, so the other readers enter without touching \code{room}, while a
  writer always needs \code{room} for itself.
  \item \textbf{Why \code{readers} needs \code{rlock}.} \code{++readers} is the same three-step
  operation as \code{counter++} in Lab 11. Without the lock two readers could both see the
  value 1, and the second would block forever on \code{room}; or an update could be lost, and
  then the room is never released, or it is released while a reader is still inside.
  \item \textbf{Readers that keep arriving.} The counter never returns to 0, so nobody posts
  \code{room} and writer 2 waits forever, although nothing is broken: starvation.
\end{enumerate}

\section{Part 4. Build a deadlock, then remove it}

Two threads, two mutexes, taken in opposite order. In the original version \code{two()} locks B
first and then A:

\begin{lstlisting}[style=ccode,caption={two() in the original deadlock.c}]
void *two(void *arg) {
    pthread_mutex_lock(&B);
    printf("two: I have B, I want A\n");
    sleep(1);
    pthread_mutex_lock(&A);
    printf("two: I have both\n");
    pthread_mutex_unlock(&A);
    pthread_mutex_unlock(&B);
    return NULL;
}
\end{lstlisting}

\begin{lstlisting}[caption={The deadlock}]
$ gcc -Wall -pthread -o deadlock deadlock.c
$ timeout 5 ./deadlock; echo "exit code: $?"
one: I have A, I want B
two: I have B, I want A
exit code: 124
\end{lstlisting}

The fix: \code{two()} locks A first, then B, the same order as \code{one()}. This is the
version handed in.

\begin{lstlisting}[style=ccode,caption={deadlock.c, fixed}]
#include <stdio.h>
#include <unistd.h>
#include <pthread.h>

pthread_mutex_t A = PTHREAD_MUTEX_INITIALIZER;
pthread_mutex_t B = PTHREAD_MUTEX_INITIALIZER;

void *one(void *arg) {
    pthread_mutex_lock(&A);
    printf("one: I have A, I want B\n");
    sleep(1);
    pthread_mutex_lock(&B);
    printf("one: I have both\n");
    pthread_mutex_unlock(&B);
    pthread_mutex_unlock(&A);
    return NULL;
}

void *two(void *arg) {
    pthread_mutex_lock(&A);
    printf("two: I have B, I want A\n");
    sleep(1);
    pthread_mutex_lock(&B);
    printf("two: I have both\n");
    pthread_mutex_unlock(&B);
    pthread_mutex_unlock(&A);
    return NULL;
}

int main(void) {
    pthread_t t1, t2;
    pthread_create(&t1, NULL, one, NULL);
    pthread_create(&t2, NULL, two, NULL);
    pthread_join(t1, NULL);
    pthread_join(t2, NULL);
    printf("done\n");
    return 0;
}
\end{lstlisting}

\begin{lstlisting}[caption={The fixed version}]
$ gcc -Wall -pthread -o deadlock deadlock.c
$ timeout 5 ./deadlock; echo "exit code: $?"
two: I have B, I want A
two: I have both
one: I have A, I want B
one: I have both
done
exit code: 0
\end{lstlisting}

In this run \code{two} got A first, so it finished before \code{one}; its first message was
left as in the original, although it now holds A at that point.

\begin{lstlisting}[caption={The original version without the two sleep(1) lines, ten runs}]
$ for i in 1 2 3 4 5 6 7 8 9 10; do timeout 2 ./deadlock > /dev/null; echo -n "$? "; done; echo
124 0 124 124 124 124 124 124 124 124
\end{lstlisting}

\textbf{Observations:}
\begin{enumerate}
  \item \textbf{The circle.} \code{one} holds A and waits for B; \code{two} holds B and waits
  for A: one \(\rightarrow\) B \(\rightarrow\) two \(\rightarrow\) A \(\rightarrow\) one. Nobody
  can go on, because each waits for a mutex that only the other can release.
  \item \textbf{Without the sleeps.} It no longer always hangs: on my laptop nine of ten runs
  hung and one finished, because sometimes one thread takes both mutexes before the other
  starts. That is worse, not better: the bug is still there but shows up only sometimes, so a
  test can pass and the program can still hang later.
  \item \textbf{Same order.} If every thread takes A before B, a thread that waits for B
  already holds A, and the owner of B never needs to wait for A, so a circle cannot form.
\end{enumerate}

\plainsection{Conclusions}

This lab combined the tools of Lab 11 into complete solutions. Producer and consumer need three
of them at once: one semaphore counts the free slots and stops a fast producer, another counts
the filled slots and stops the consumer on an empty buffer, and a mutex protects the buffer.
The condition variable solves the same waiting problem in a different style: the state is kept
in an ordinary variable under a mutex, and the condition variable is only a place to sleep,
which is why the condition is checked again in a loop after waking. Readers and writers showed
that a lock can be shared: the first reader takes the room and the last one returns it.

The lab also showed what goes wrong when the tools are combined carelessly. Sleeping while
holding a mutex blocked producer and consumer after four items, and two mutexes taken in
opposite order stopped two threads forever. Both are deadlocks, both come from the order of the
operations, and both disappear when the order is fixed: wait first and lock second, and take
locks in the same order everywhere. A correct program can still be unfair, as the starving
writer shows.

\end{document}
